% ==========================================
% PARAM.TEX - Configurazione Ingegneria del Software
% ==========================================

% --- 1. LINGUA E FONT ---
\usepackage[T1]{fontenc}
\usepackage[utf8]{inputenc}
\usepackage[italian]{babel} % Traduce "Chapter" in "Capitolo", ecc.
\usepackage{lmodern}        % Font vettoriale di alta qualità

% --- 2. GEOMETRIA E IMPAGINAZIONE ---
\usepackage[a4paper, top=2.5cm, bottom=2.5cm, left=2.5cm, right=2.5cm]{geometry}
\usepackage{titlesec}       % Personalizzazione dei titoli

% Rimuove la scritta "Capitolo X" e affianca il numero al titolo
\titleformat{\chapter}[hang]
  {\normalfont\huge\bfseries}{\thechapter.}{20pt}{\Huge}
\titlespacing*{\chapter}{0pt}{-20pt}{30pt}

% --- 3. MATEMATICA E SIMBOLI ---
\usepackage{amsmath, amssymb, amsfonts}

% --- 4. GRAFICA E IMMAGINI ---
\usepackage{graphicx}
\usepackage{float}          % Permette l'uso di [H] per forzare la posizione delle immagini
\usepackage{caption}        % Personalizzazione delle didascalie
\usepackage{subcaption}     % Per affiancare più immagini

% --- 5. TIKZ E UML ---
\usepackage{tikz}
% NOTA: Il file tikz-uml.sty deve essere nella stessa cartella del progetto
\usepackage{tikz-uml} 

% --- 6. LISTE E TABELLE ---
\usepackage{enumitem}       % Controllo avanzato di itemize e enumerate
\usepackage{booktabs}       % Tabelle professionali (toprule, midrule, bottomrule)
\usepackage{tabularx}       % Tabelle che si adattano alla larghezza della pagina

% --- 7. FORMATTAZIONE DEL CODICE (LISTINGS) ---
\usepackage{listings}
\usepackage{xcolor}

% Definizione colori per il codice
\definecolor{codegreen}{rgb}{0,0.6,0}
\definecolor{codegray}{rgb}{0.5,0.5,0.5}
\definecolor{codepurple}{rgb}{0.58,0,0.82}
\definecolor{backcolour}{rgb}{0.95,0.95,0.92}

% Configurazione stile codice (ottimizzato per Java/C++)
\lstdefinestyle{oopstyle}{
    backgroundcolor=\color{backcolour},   
    commentstyle=\color{codegreen}\itshape,
    keywordstyle=\color{magenta}\bfseries,
    numberstyle=\tiny\color{codegray},
    stringstyle=\color{codepurple},
    basicstyle=\ttfamily\footnotesize,
    breakatwhitespace=false,         
    breaklines=true,                 
    captionpos=b,                    
    keepspaces=true,                 
    numbers=left,                    
    numbersep=5pt,                  
    showspaces=false,                
    showstringspaces=false,
    showtabs=false,                  
    tabsize=2,
    frame=single,
    rulecolor=\color{black!30}
}
\lstset{style=oopstyle}

% --- 8. BOX COLORATI (PATTERN GRASP, GOF, PROCESSO UP) ---
\usepackage[most]{tcolorbox}

% Box per i Pattern GRASP (Azzurro)
\newtcolorbox{graspbox}[1]{
    enhanced, breakable, 
    colback=blue!5!white, colframe=blue!75!black,
    title=\textbf{Pattern GRASP:} #1,
    fonttitle=\bfseries,
    boxrule=1pt, arc=3pt, drop shadow
}

% Box per i Design Pattern GoF (Verde)
\newtcolorbox{patternbox}[1]{
    enhanced, breakable, 
    colback=green!5!white, colframe=green!50!black,
    title=\textbf{Design Pattern:} #1,
    fonttitle=\bfseries,
    boxrule=1pt, arc=3pt, drop shadow
}

% Box per le Fasi del Processo UP (Arancione)
\newtcolorbox{upbox}[1]{
    enhanced, breakable, 
    colback=orange!5!white, colframe=orange!80!black,
    title=\textbf{Fase UP:} #1,
    fonttitle=\bfseries,
    boxrule=1pt, arc=3pt, drop shadow
}

% --- 9. COLLEGAMENTI E INDICE (Da caricare per ultimo) ---
\usepackage{hyperref}
\hypersetup{
    colorlinks=true,
    linkcolor=blue!70!black, % Colore dei link interni (indice)
    filecolor=magenta,       % Colore dei link a file
    urlcolor=cyan,           % Colore degli URL
    pdftitle={Appunti di Analisi e Progettazione del Software},
    pdfauthor={Il tuo nome},
    pdfpagemode=UseOutlines
}
\usepackage{bookmark} % Migliora i segnalibri nel lettore PDF